
%%FILE blueprint/src/chapter/Almost.tex
\section{Almost Mathematics}

%%FILE blueprint/src/chapter/Cyclotomic.tex
\section{Calculation of $\Q_p^{\cyc}$}

%%FILE blueprint/src/chapter/Infrastructure.tex
\section{Missing Infrastructures in Mathlib}
In this section, we record the missing infrastructures in a rough way. All definitions and theorem under a same topic will be bundled in a single definition.

\begin{definition}
    \label{Relate Galois Group with Category of Finite Field Extensions}
    \notready

    Inteperate Galois group as functors on $ \text{FinExt}_K $ satisfying certain property. Show that $\text{FinExt}_K \simeq \text{FinExt}_K'$ induces an isomorphism $\Gal_K \simeq \Gal_{K'}$.
\end{definition}

%%FILE blueprint/src/chapter/PerfectoidField.tex
\section{Perfectoid Fields}

\begin{definition}[Perfectoid Fields]
    \label{Perfectoid Field}
    \lean{PerfectoidField}
    % \uses{Valuation.integer}

    Let $K$ be a topological field. We say that $K$ is a \emph{perfectoid field} if and only if the following conditions are satisfied:\\
    (a) the topology on K is induced by a valuation $\| \cdot \| : K \to \R_{\ge0}$, and $K$ is complete for this topology;\\
    (b) there exists a non-unit $\pi \in \O_K$ such that $p \in \pi^p \O_K$;\\
    (c) every element of $\O_K/p\O_K$ is a $p$th-power.
\end{definition}

%%FILE blueprint/src/chapter/TiltingPerfectoidField.tex

\section{The Tilting Equivalence for Perfectoid Fields}

\begin{definition}
    \label{Tilting Functor}
    \notready

    There is an equivalence of category
    \begin{align*}
        (-)^\flat : \left \{ \text{Perfectoid fields over } K \right \} & \to \left \{ \text{Perfectoid fields over } K^\flat \right \} : (-)^\# \\
        L & \mapsto L^\flat
    \end{align*}

\end{definition}

\begin{theorem}
    \label{untilt_same_degree}
    \uses{Tilting Functor}
    \notready

    Let $M$ be a finite extension of $K^\flat$. ($M$ is automatically a perfectoid field. add this theorem) Then its untilt $M^\#$ is of the same degree. 

\end{theorem}

\begin{proof}
    \notready

    sorry
\end{proof}



\begin{definition}
    \label{Tilting Functor (Finite Extension)}
    \uses{Tilting Fuctor}
    \notready

    There is a functor 
    \begin{align*}
        \left \{ \text{finite field extensions of } K \right \} & \to \left \{ \text{finite fields extension of} K^\flat \right \} \\
        L & \mapsto L^\flat
    \end{align*}

\end{definition}

\begin{theorem}
    \label{tilting_finite_degree_preserving}
    \uses{Tilting Functor (Finite Extension)}
    \notready

    The functor above is degree preserving.
\end{theorem}

\begin{proof}
    \notready

    sorry
\end{proof}

\begin{theorem}
    \label{tilting_finite_equivalence_of_category}
    \uses{Tilting Functor (Finite Extension)}
    \notready

    The functor above is an equivalence of category.
\end{theorem}

\begin{proof}
    \uses{tilting_finite_degree_preserving}
    \notready

    sorry
\end{proof}

\begin{definition}[Galois Isomorphism of Perfectoid Fields]
    \label{Galois Isomorphism of Perfectoid Fields}

    \uses{
        Tilting Functor (Finite Extension), 
        tilting_finite_equivalence_of_category, 
        Relate Galois Group with Category of Finite Field Extensions
        }
    \notready

    There exists an isomorphism of absolute Galois groups
    $$ \Gal_K \simeq \Gal_{K^\flat}.$$

    \emph{Do we need write down the explicit isomorphism?}
\end{definition}


%%FILE blueprint/src/content.tex
\title{Fontaine-Wintenberger Theorem}


% \home{https://lean4fonwinteam.github.io/FontaineWintenberger/}
%\home{localhost:8080}
\input{homeurl}
\github{https://github.com/Lean4FonWinTeam/FontaineWintenberger/}
%\dochome{https://lean4fonwinteam.github.io/FontaineWintenberger/doc/}
%\dochome{localhost:8080/doc}

\maketitle


\tableofcontents

The first section records currently missing infrastructures in Mathlib. The following three sections serves for the tilting equivalence. The remaining parts serves for the concrete example of cyclotomic fields.

\input{\MyPath/chapter/Infrastructure.tex}

\input{\MyPath/chapter/PerfectoidField.tex}

\input{\MyPath/chapter/Almost.tex}

\input{\MyPath/chapter/TiltingPerfectoidField.tex}

\input{\MyPath/chapter/Cyclotomic.tex}

\section{Example}
Here you can use LaTex. \cite{marcus}. You must cite something to make it work.

%\input{demo}


%%FILE blueprint/src/demo.tex

\begin{lemma}\label{List.aux}
        \leanok
        \lean{List.aux}
                this should be a doc string.
    \end{lemma}

\begin{proof}
    \leanok
\end{proof}

\begin{theorem}\label{List.Nat.Ex}
        \leanok
        \lean{List.Nat.Ex}
        \uses{List.aux}
                Here is a doc string
    \end{theorem}

\begin{proof}
    \leanok
\end{proof}

\begin{lemma}\label{List.test}
        \lean{List.test}
                No documentation.
    \end{lemma}

\begin{theorem}\label{List.Test}
        \lean{List.Test}
        \uses{List.test}
                No documentation.
    \end{theorem}



%%FILE blueprint/src/homeurl.tex
\home{localhost:8080}
\dochome{localhost:8080/doc}

%%FILE blueprint/template.tex

((* for full_name, features in structure_dict.statements.items() *))
((* if features[1] in ["theorem", "lemma", "definition"] *))
\begin{((( features[1] )))}\label{((( full_name )))}
    ((* if features[2] *))
    \leanok
    ((* endif *))
    \lean{((( full_name )))}
    ((* if structure_dict.dependencies[full_name] *))
    \uses{((( ','.join(structure_dict.dependencies[full_name]) )))}
    ((* endif *))
    ((* if structure_dict.docs[full_name] *))
        ((( structure_dict.docs[full_name] )))
    ((* else *))
        No documentation.
    ((* endif *))
\end{((( features[1] )))}

((* endif *))
((* if features[1] in ["theorem", "lemma"] and features[2] *))
\begin{proof}
    \leanok
\end{proof}

((* endif *))
((* endfor *))
